\paragraph{冻结尾部的凸延拓刚性、Legendre 硬壁与 Abel 有限部余项}
在定理 \ref{thm:xi-time-part57b-pressure-spectrum-discrete-convexity-monotone-excess} 已把整数压力谱锚定为离散凸序列、定理 \ref{thm:xi-freezing-affine-rigidity-two-point-recovery} 已把冻结相整数尾部压缩为两点即可恢复的仿射律之后，还可继续把连续凸延拓、Legendre 对偶边界与 Möbius--Euler 截断余项进一步收束为下列四条后果。

\begin{theorem}[冻结阈值之后任意凸压力延拓的仿射刚性]\label{thm:xi-time-part57f-convex-extension-affine-tail}
\leanverified{paper\_xi\_time\_part57f\_convex\_extension\_affine\_tail}
在严格基态间隙
\[
\alpha_*^{(2)}<\alpha_*
\]
下，记
\[
a_{\mathrm c}:=\frac{\log\varphi-g_*}{\alpha_*-\alpha_*^{(2)}},
\qquad
a_0:=\lfloor a_{\mathrm c}\rfloor+1.
\]
设
\[
\Lambda:[a_0,\infty)\to\RR
\]
为任意凸函数，并满足
\[
\Lambda(n)=P_n
\qquad
\bigl(n\in\ZZ,\ n\ge a_0\bigr).
\]
则必有
\[
\boxed{\;
\Lambda(t)=\alpha_* t+g_*
\qquad
(t\ge a_0).\;}
\]
特别地，任何保持凸性并在冻结区整数锚点与离散压力序列对齐的连续压力外推，在阈值之后都不再允许残留曲率；因此，定理 \ref{thm:xi-projective-pressure-canonical-convex-majorant} 所给出的连续压力接口一旦在这些整数锚点上实现精确对齐，其尾部也只能是同一条仿射直线。
\end{theorem}
\begin{proof}
由定理 \ref{thm:fold-pressure-freezing-threshold}，对每个整数 \(n\ge a_0\) 都有
\[
P_n=n\alpha_*+g_*,
\]
故
\[
G(t):=\Lambda(t)-(\alpha_* t+g_*)
\]
是定义在 \([a_0,\infty)\) 上的凸函数，并满足
\[
G(n)=0
\qquad
\bigl(n\in\ZZ,\ n\ge a_0\bigr).
\]

先看首个区间 \([a_0,a_0+1]\)。端点 \(t=a_0,a_0+1\) 时结论显然，故只需取
\[
t\in(a_0,a_0+1).
\]
由凸函数割线斜率单调性，对
\[
t<a_0+1<a_0+2
\]
应用
\[
\frac{G(a_0+1)-G(t)}{a_0+1-t}
\le
\frac{G(a_0+2)-G(a_0+1)}{1},
\]
便得
\[
\frac{-G(t)}{a_0+1-t}\le 0,
\]
从而 \(G(t)\ge 0\)。另一方面，由
\[
a_0\le t\le a_0+1
\]
与凸性，
\[
G(t)\le
\frac{a_0+1-t}{1}G(a_0)+\frac{t-a_0}{1}G(a_0+1)=0,
\]
故 \(G(t)=0\)。于是
\[
G\equiv 0
\qquad\text{于 }[a_0,a_0+1].
\]

再取任意整数 \(n\ge a_0+1\)。端点 \(t=n,n+1\) 时同样平凡，故只需考虑
\[
t\in(n,n+1).
\]
同样由割线斜率单调性，对
\[
n-1<n\le t
\]
有
\[
\frac{G(n)-G(n-1)}{1}
\le
\frac{G(t)-G(n)}{t-n},
\]
即
\[
0\le \frac{G(t)}{t-n},
\]
故 \(G(t)\ge 0\)。而由 \(G(n)=G(n+1)=0\) 与凸性，
\[
G(t)\le
\frac{n+1-t}{1}G(n)+\frac{t-n}{1}G(n+1)=0.
\]
因此 \(G(t)=0\)。这表明 \(G\equiv 0\) 于每个区间 \([n,n+1]\) 上成立。并合全部相邻整数区间，即得
\[
G(t)\equiv 0
\qquad
(t\ge a_0),
\]
亦即
\[
\Lambda(t)=\alpha_* t+g_*.
\]
证毕。
\end{proof}

\begin{theorem}[冻结尾部的 Legendre 硬壁与边界代价]\label{thm:xi-time-part57f-legendre-hard-wall-boundary-cost}
沿用定理 \ref{thm:xi-time-part57f-convex-extension-affine-tail} 的记号。在 \(X_m\) 上取均匀分布 \(u_m\)，并定义随机变量
\[
Y_m(x):=\frac{1}{m}\log d_m(x)
\qquad
(x\in X_m).
\]
则对每个整数 \(a\ge 1\)，有
\[
\frac1m\log \EE_{u_m}\!\left[e^{amY_m}\right]
=
\frac1m\log\frac{S_a(m)}{|X_m|}
\longrightarrow
P_a-\log\varphi.
\]
再设
\[
I(y):=
\sup_{t\ge a_0}
\Bigl\{
ty-\bigl(\Lambda(t)-\log\varphi\bigr)
\Bigr\}
\]
为任一此类凸延拓 \(\Lambda\) 所对应的 Legendre--Fenchel 对偶。则成立：
\begin{enumerate}
  \item 对任意 \(y>\alpha_*\)，有
  \[
  \boxed{\;I(y)=+\infty.\;}
  \]
  \item 在边界点 \(y=\alpha_*\) 处，恰有
  \[
  \boxed{\;I(\alpha_*)=\log\varphi-g_*.\;}
  \]
\end{enumerate}
换言之，\(\alpha_*\) 是冻结尾部在 Legendre 对偶中的硬上壁，而抵达该边界的指数代价恰由环境熵率与基态简并指数之差给出。
\end{theorem}
\begin{proof}
由定义
\[
\EE_{u_m}\!\left[e^{amY_m}\right]
=
\frac{1}{|X_m|}
\sum_{x\in X_m} d_m(x)^a
=
\frac{S_a(m)}{|X_m|}.
\]
又因 \(|X_m|=F_{m+2}\) 且
\[
\lim_{m\to\infty}\frac1m\log F_{m+2}=\log\varphi,
\]
故
\[
\frac1m\log \EE_{u_m}\!\left[e^{amY_m}\right]
=
\frac1m\log S_a(m)-\frac1m\log|X_m|
\longrightarrow
P_a-\log\varphi.
\]

另一方面，由定理 \ref{thm:xi-time-part57f-convex-extension-affine-tail}，
\[
\Lambda(t)=\alpha_* t+g_*
\qquad
(t\ge a_0).
\]
于是对任意 \(t\ge a_0\)，都有
\[
ty-\bigl(\Lambda(t)-\log\varphi\bigr)
=
t(y-\alpha_*)+\log\varphi-g_*.
\]
若 \(y>\alpha_*\)，右端随 \(t\to+\infty\) 线性发散到 \(+\infty\)，从而
\[
I(y)=+\infty.
\]
若 \(y=\alpha_*\)，则上式对所有 \(t\ge a_0\) 都恒等于
\[
\log\varphi-g_*,
\]
故其上确界亦为该常数，即
\[
I(\alpha_*)=\log\varphi-g_*.
\]
证毕。
\end{proof}

\begin{theorem}[Abel 有限部常数的内禀 \texorpdfstring{$\eta$}{eta}-截断误差界]\label{thm:xi-time-part57f-abel-finite-part-eta-tail}
沿用命题 \ref{prop:real-input-40-logM-infty} 的记号。记
\[
H(x):=(1-x)(1-6x+9x^2-x^3-x^4)
=
1-7x+15x^2-10x^3+x^5.
\]
并定义
\[
\eta:=\min_{0\le x\le b^2}H(x)>0.
\]
对每个整数 \(N\ge 2\)，记
\[
\log\mathfrak M_\infty^{(N)}
:=
\log C_\infty
-\sum_{m=2}^{N}\frac{\mu(m)}{m}\log H(b^m).
\]
则截断误差满足显式上界
\[
\boxed{\;
\bigl|\log\mathfrak M_\infty-\log\mathfrak M_\infty^{(N)}\bigr|
\le
\frac{7}{\eta(1-b)}\cdot \frac{b^{N+1}}{N+1}.\;}
\]
因此，Abel 有限部常数的 Möbius--Euler 展开不只几何收敛，而且其余项可由区间 \([0,b^2]\) 上的内禀最小值 \(\eta\) 直接审计。
\end{theorem}
\begin{proof}
由于 \(b\) 是多项式
\[
1-6x+9x^2-x^3-x^4
\]
的最大正根，故对 \(0\le x\le b^2<b\) 有
\[
1-6x+9x^2-x^3-x^4>0
\qquad\text{且}\qquad
1-x>0.
\]
因此
\[
H(x)>0
\qquad
(0\le x\le b^2),
\]
从而 \(\eta>0\)。

又由命题 \ref{prop:real-input-40-logM-infty} 的同一根定位，可知
\[
\frac14<b<\frac13,
\]
于是
\[
b^m\le b^2<\frac19
\qquad
(m\ge 2).
\]
对 \(x\in[0,b^2]\) 直接计算得
\[
1-H(x)=7x-15x^2+10x^3-x^5\le 7x.
\]
另一方面，由 \(x\le 1/9\) 可知
\[
1-H(x)
=
x\bigl(7-15x+10x^2-x^4\bigr)
\ge
x\Bigl(7-\frac{15}{9}\Bigr)\ge 0,
\]
故
\[
0<H(x)\le 1
\qquad
(0\le x\le b^2).
\]
因此对每个 \(m\ge 2\)，都可应用基本不等式
\[
-\log t\le \frac{1-t}{t}
\qquad
(0<t\le 1),
\]
得到
\[
\bigl|\log H(b^m)\bigr|
=
-\log H(b^m)
\le
\frac{1-H(b^m)}{H(b^m)}
\le
\frac{7b^m}{\eta}.
\]
于是
\[
\bigl|\log\mathfrak M_\infty-\log\mathfrak M_\infty^{(N)}\bigr|
\le
\sum_{m>N}\frac{|\mu(m)|}{m}\,\bigl|\log H(b^m)\bigr|
\le
\frac{7}{\eta}\sum_{m>N}\frac{b^m}{m}.
\]
而对 \(m>N\) 有 \(1/m\le 1/(N+1)\)，故
\[
\sum_{m>N}\frac{b^m}{m}
\le
\frac{1}{N+1}\sum_{m>N}b^m
=
\frac{1}{N+1}\cdot \frac{b^{N+1}}{1-b}.
\]
代回即得所示盒装上界。证毕。
\end{proof}

\begin{proposition}[Möbius--Euler 余项的首项由下一平方自由层主导]\label{prop:xi-time-part57f-mobius-euler-next-squarefree-layer}
沿用定理 \ref{thm:xi-time-part57f-abel-finite-part-eta-tail} 的记号。定义截断余项
\[
R_N:=
\log\mathfrak M_\infty-\log\mathfrak M_\infty^{(N)}
=
-\sum_{m>N}\frac{\mu(m)}{m}\log H(b^m).
\]
则当 \(N\to\infty\) 时，有统一渐近展开
\[
\boxed{\;
R_N
=
7\sum_{m>N}\frac{\mu(m)}{m}b^m
+
O\!\left(\frac{b^{2N+2}}{N+1}\right).\;}
\]
特别地，若 \(N+1\) 为平方自由数，则
\[
\boxed{\;
R_N
=
\frac{7\,\mu(N+1)}{N+1}b^{N+1}
+
O\!\left(\frac{b^{N+2}}{N+2}\right).\;}
\]
因此，余项的首个非零主项由下一层平方自由指数唯一决定，其符号直接由 \(\mu(N+1)\) 控制。
\end{proposition}
\begin{proof}
由
\[
H(x)=1-7x+O(x^2)
\qquad
(x\to 0),
\]
可得
\[
\log H(x)=-7x+O(x^2)
\qquad
(x\to 0).
\]
把 \(x=b^m\) 代入并带回余项展开，得到
\[
\begin{aligned}
R_N
&=
-\sum_{m>N}\frac{\mu(m)}{m}\Bigl(-7b^m+O(b^{2m})\Bigr)\\
&=
7\sum_{m>N}\frac{\mu(m)}{m}b^m
+
O\!\left(\sum_{m>N}\frac{b^{2m}}{m}\right).
\end{aligned}
\]
又因
\[
\sum_{m>N}\frac{b^{2m}}{m}
\le
\frac{1}{N+1}\sum_{m>N}b^{2m}
=
O\!\left(\frac{b^{2N+2}}{N+1}\right),
\]
便得第一条渐近式。

若 \(N+1\) 为平方自由数，则 \(\mu(N+1)\neq 0\)，从而
\[
\sum_{m>N}\frac{\mu(m)}{m}b^m
=
\frac{\mu(N+1)}{N+1}b^{N+1}
+
O\!\left(\sum_{m\ge N+2}\frac{b^m}{m}\right).
\]
而
\[
\sum_{m\ge N+2}\frac{b^m}{m}
\le
\frac{1}{N+2}\sum_{m\ge N+2}b^m
=
O\!\left(\frac{b^{N+2}}{N+2}\right).
\]
将其代回第一条式子，即得
\[
R_N
=
\frac{7\,\mu(N+1)}{N+1}b^{N+1}
+
O\!\left(\frac{b^{N+2}}{N+2}\right).
\]
证毕。
\end{proof}

\noindent 结合定理 \ref{thm:xi-time-part57b-pressure-spectrum-discrete-convexity-monotone-excess} 的离散凸性、定理 \ref{thm:xi-time-part57f-convex-extension-affine-tail} 的凸延拓刚性、定理 \ref{thm:xi-time-part57f-legendre-hard-wall-boundary-cost} 的 Legendre 边界代价，以及定理 \ref{thm:xi-time-part57f-abel-finite-part-eta-tail} 与命题 \ref{prop:xi-time-part57f-mobius-euler-next-squarefree-layer} 对 Möbius--Euler 余项的显式控制可见，冻结相在压力对偶与 Abel 有限部两侧都呈现出同一种硬壁结构：一侧把可达纤维指数锁定在 \(y\le \alpha_*\) 并在边界处留下精确代价 \(\log\varphi-g_*\)，另一侧则把有限部余项压缩为由下一平方自由层主导的首项与可闭式审计的尾界。这为进一步把谱硬壁与有限部硬余项统一到同一 Legendre--Witt 框架提供了直接接口。

\endinput
